% ===========================================================================
%  utente/03-studio/00-capitolo.tex
% ===========================================================================
\chapter{Studio}\label{cap:studio}

Lo Studio è dove si costruiscono le cose. Ci si arriva da \voce{Studio} nel
menu, oppure dal pulsante \pulsante{Open Studio} nella schermata iniziale.

\figuraSchermata{1.0}{figure/schermate/studio.png}
  {Lo Studio. In alto le tre modalità di lavoro, sotto le cinque famiglie di
   artefatti, al centro il tavolo da disegno.}
  {studio}

\section{Le cinque famiglie}\label{sez:studio-famiglie}

La riga di schede sotto il titolo dice \emph{che cosa} si sta costruendo.
Sono cinque cose diverse, e vale la pena sapere quale serve prima di
cominciare.

\begin{center}
\begin{tabular}{@{}L{28mm}L{110mm}@{}}
\toprule
\textbf{Scheda} & \textbf{Che cos'è} \\
\midrule
\rigaTenue
\voce{Strategies} & Una strategia: decide che cosa comprare e vendere, e
  quando. È la cosa che si costruisce più spesso. \\
\voce{Approaches} & Un approccio: salta tutto il motore e produce
  direttamente la serie dei guadagni e delle perdite. Serve a chi ha già un
  proprio calcolo e vuole solo farlo entrare nel sistema. \\
\rigaTenue
\voce{Weight mode} & Un criterio per distribuire il capitale fra più
  strategie. Ne sono già installati quattordici. \\
\voce{Overlay} & Un freno: riduce l'esposizione quando il mercato si
  comporta in un certo modo. Non guadagna, protegge. \\
\rigaTenue
\voce{Metrics} & Una misura da aggiungere a quelle già presenti nelle
  schede dei risultati. \\
\bottomrule
\end{tabular}
\end{center}

\section{I tre modi di costruire}\label{sez:studio-modi}

In alto a destra si sceglie come lavorare. Il risultato è lo stesso --- un
file di codice installato nel sistema --- e cambia solo la strada.

\subsection{Blueprint: a blocchi}

Si trascinano blocchi dalla libreria di sinistra sul tavolo e si collegano le
uscite agli ingressi. La libreria è divisa per categorie: prezzi, numeri,
dati esterni come utili e fondamentali.

\begin{pratica}{Il modo più rapido per capire come funziona}
Premere \pulsante{Load an example}: compare un esempio già collegato
(l'incrocio di due medie con uno stop). Lo si smonta e si rimonta: in dieci
minuti si è capito il meccanismo senza aver scritto una riga.
\end{pratica}

\subsection{Code: scrivendo}

Il pulsante \pulsantesec{\{ \} Code} passa all'editor di testo. È lo stesso
artefatto, visto come codice Python. Si può passare da una vista all'altra.

\subsection{AI agent: a parole}

La scheda \voce{AI agent} apre una conversazione: si descrive a parole la
strategia e il modello la scrive. Richiede un modello installato in locale;
senza, la scheda lo dice e le altre due modalità continuano a funzionare.

\figuraSchermata{1.0}{figure/schermate/agent.png}
  {L'assistente. Si descrive che cosa si vuole, e lui scrive l'artefatto e lo
   installa.}
  {agent}

\begin{attenzione}{Quello che scrive l'assistente va letto}
Il codice generato viene installato nel sistema e gira sui propri dati --- e
un giorno, forse, sul proprio conto. L'assistente non sa se una strategia ha
senso dal punto di vista finanziario: sa scrivere del codice che fa quello
che gli è stato chiesto. La scheda \voce{Code} serve anche a questo.
\end{attenzione}

\section{Genetics}\label{sez:studio-genetica}

La scheda \voce{Genetics} apre una zona a sé: invece di scrivere una
strategia, se ne fanno nascere molte e si tengono quelle che sopravvivono.
Si sceglie l'universo di strumenti e si lascia correre.

\figuraSchermata{1.0}{figure/schermate/genetic.png}
  {La sezione genetica. Si dichiara su che cosa far evolvere le strategie e
   il resto è attesa.}
  {genetic}

\begin{nota}{Dura a lungo}
È il lavoro più lento dell'applicazione. Parte come processo in sottofondo:
si può chiudere la pagina, e la campanella avvisa a conclusione.
\end{nota}

\section{Che cosa succede quando si preme «Install»}\label{sez:studio-install}

L'artefatto viene scritto in una cartella del sistema e da quel momento
compare negli elenchi: una strategia fra quelle lanciabili, una modalità di
pesatura fra quelle scegliibili in un project, un overlay fra quelli
attivabili.

Prima di scrivere, il sistema controlla il codice: non deve leggere o
scrivere file, non deve aprire connessioni, non deve eseguire comandi. Se
contiene una di queste cose l'installazione viene rifiutata e il motivo è
scritto.
